\documentclass[../../main.tex]{subfiles}

\begin{document}

\subsection{Il calorimetro elettromagnetico}
  \label{subsec:il-calorimetro-elettromagnetico}
  \begin{figure}[ht]
    \centering
    \includegraphics[width=.5\textwidth]{DUNE/sand/ecal}
    \caption{
      Immagine del calorimetro elettromagnetico dell'esperimento KLOE,
      in cui sono visibili sia i moduli del barile (\textit{barrel})
      che quelli dei tappi (\textit{endcaps}).~\cite{Pia_2024}
    }
    \label{fig:ecal}
  \end{figure}
  Il calorimetro elettromagnetico (ECAL), schematizzato in figura
  ~\ref{fig:cal-mag} e fotografato in figura~\ref{fig:ecal}, è un calorimetro a
  campionamento a piombo e fibre scintillanti.
  Le fibre scintillanti garantiscono un'efficiente trasmissione della luce su
  diversi metri, offrono una risoluzione temporale sub-nanometrica e un'ottima
  ermeticità.

  Il rilevatore è composto da: un calorimetro a barile, di forma quasi
  cilindrica, situato all'interno del magnete vicino al
  criostato della bobina; e due calorimetri aggiuntivi, chiamati tappi
  (\textit{end-caps}), posizionati alle estremità del magnete per garantire
  l'ermeticità.
  Il barile è composto da 24 moduli, ciascuno dei quali misura 4.3~m di
  lunghezza, 23~cm di spessore e presenta una sezione trasversale
  trapezoidale con basi larghe 52~cm e 59~cm.
  I tappi sono costituiti da 32 moduli rettangolari verticali, spessi 23~cm
  e con lunghezze variabili da 0,7~m a 3,9~m.
  Questi moduli sono piegati alle estremità superiore e inferiore per consentire
  l'inserimento nel calorimetro a barile e per
  posizionare gli assi dei fototubi parallelamente al campo magnetico.\\
  Grazie all'ampia sovrapposizione tra il barile e i tappi, il calorimetro
  ottiene una copertura completa poiché non presenta gap inattivi
  all'interfaccia tra questi componenti.

  Il calorimetro dispone di risoluzione spaziale di 1,3~cm nelle coordinate
  $r - \phi$ per il barile e nelle coordinate $x - z$ per i tappi.
  Le risoluzioni in energia e tempo sono state valutate durante le fasi
  operative di KLOE~\cite{ADINOLFI2002364} e sono:
  \begin{align*}
    \frac{\sigma_E}{E} = \frac{5 \%}{\sqrt{E}} &&
    \sigma_t = \frac{\SI{54}{ps}}{\sqrt{E}} && \text{(Energie in GeV)}
  \end{align*}


\end{document}